% !TeX root = ../../../main.tex
\begin{abstract}
Quantum computing is often introduced from the inside out: first qubits, then gates, then an application to justify the machine. This perspective and framework paper reverses that order, presenting a mechanism-first, resource-bounded method for deciding whether a research question suits a quantum solution and, if so, how to engineer the algorithm around a physical device rather than an imaginary one. The method formalizes the problem, its input access, output contract, error tolerance, verification path, and classical baseline before selecting a quantum primitive. It then develops end-to-end resource calculations for noisy and fault-tolerant hardware, from shot budgets and time-to-solution to logical error budgets, physical-qubit overheads, and classical feedback latency. A hybrid partition is treated as an optimization problem in its own right: the quantum processor receives only the coherent, reusable kernel whose benefit exceeds the cost of crossing the interface.

The framework is then used to audit proposed framings in DNA-style phase encodings, constrained scheduling, and lattice gauge theory; a small, fully reproducible execution of the whole audit---a four-qubit Ising instance checked against an independent classical oracle and run once, uncorrected, on real hardware---is included as an appendix. The audits expose a consistent lesson: native quantum dynamics, compact observables, exploitable symmetry, reversible predicates, and structured spectral questions illuminate real opportunity; large data, generic hardness, expensive loading, full-vector outputs, dense embeddings, and unverified heuristics usually do not. The paper proves no new theorem and claims no quantum speedup or advantage. The result is not a catalog of quantum algorithms but a disciplined way to ask whether the quantum work is actually there.
\end{abstract}

\noindent\textbf{Keywords:} \PaperKeywords
